\subsubsection{Initialization versus adversarial training}
\label{sec:initialization-ablation}

The lognormal-shaped initialization could itself account for the preceding
pilot's favorable point estimates. A retrospective paired ablation therefore
reuses all forty evaluation datasets, without new tuning or sample selection.
It holds the outer folds, inner split, numerical integration, score, and
bandwidths fixed, and replaces adversarial training by the existing data-only
initialization with zero optimizer steps. This linear normal-quantile
initialization is a lognormal-family restriction, not an unrestricted learner.
This is a mechanism diagnostic on already inspected samples, not a second
independent confirmation.

Three comparisons distinguish training from weight selection. Fixed weights
retain the fifty-pseudo-observation rule. Reselected weights allow initialized
and trained generators to choose weights separately on the same inner
validation data. Matched weights instead give the initialized outer nuisance
the trained learner's weights. The last comparison isolates outer-nuisance
training conditional on that weight-selection policy; it is not a cheap
standalone estimator, since selecting its weights still requires training.
All prior trained estimates are preserved. Thirty-two recomputed estimate
and standard-error checks agree with the archive within $10^{-15}$.

Table~\ref{tab:initialization} shows that initialization alone nearly
reproduces the lognormal adaptive results: RMSE is 0.0527 and 0.0314,
compared with 0.0525 and 0.0318 after training. Under fixed calibration, all
four training-induced RMSE changes are below 0.0004. These results do not
identify adversarial learning as the source of the earlier lognormal gains.
With separately selected weights, initialization also has lower point RMSE
in both Weibull cells, although paired Monte Carlo uncertainty remains large.

The matched-weight comparison adds an important qualification. At Weibull
$n=2000$, training reduces RMSE from 0.0585 to 0.0393 when the trained
selection rule's weights are held fixed. Its paired MSE difference is
$-0.00188$, with Monte Carlo standard error 0.00108. Yet allowing the
initialized learner to select its own weights gives RMSE 0.0366, while the
smoothed plug-in remains more accurate at 0.0328. Averaged over all folds and
distinct bandwidths in this cell, generated-component weights rise from
0.139 before training to 0.400 after training. Training and the induced
weight response can therefore move the reserve error in different directions.

Training does improve observable-distribution fit under Weibull. Mean
outer-holdout maximum-distribution $W_1$ falls from 0.0955 to 0.0780 at
$n=500$ and from 0.0657 to 0.0483 at $n=2000$, using equal integration
precision. These are descriptive means over thirty outer fits per cell;
overlapping folds are not thirty independent replications. Holdout fit is
used only for reporting. Better global fit therefore does not establish
improved reserve estimation under adaptive calibration. Figure~\ref{fig:initialization}
reports paired uncertainty: every displayed approximate Monte Carlo interval
includes zero. The ten-repetition retrospective study motivates examining
conservative mixing on fresh data, but does not validate such a modification
or justify replacing the fixed-calibration safeguard.

\begin{table}[htbp]
\centering
\linespread{1}\selectfont
\small
\caption{Exact-reserve RMSE before and after adversarial training}
\label{tab:initialization}
\resizebox{\textwidth}{!}{\begin{tabular}{lrccccc}
\toprule
 & & \multicolumn{2}{c}{Fixed weights} & \multicolumn{2}{c}{Reselected weights} & Matched \\
Design & $n$ & Init. & Trained & Init. & Trained & Init. \\
\midrule
Lognormal & 500 & 0.0577 & 0.0575 & 0.0527 & 0.0525 & 0.0521 \\
Lognormal & 2000 & 0.0335 & 0.0335 & 0.0314 & 0.0318 & 0.0316 \\
Weibull & 500 & 0.0538 & 0.0534 & 0.0486 & 0.0512 & 0.0485 \\
Weibull & 2000 & 0.0327 & 0.0327 & 0.0366 & 0.0393 & 0.0585 \\
\bottomrule
\end{tabular}
}
\begin{minipage}{0.96\textwidth}
\footnotesize
\emph{Notes:} Ten paired repetitions per cell, reusing the preceding pilot's
complete evaluation cohort. ``Init.'' has zero adversarial updates. The
matched initialized arm uses weights selected by the trained learner, so its
comparator is the trained column under reselected weights. Fixed-bandwidth
results, coverage, all raw estimates, validation risks, and cache provenance
are included in the replication archive. No root failures occur in this ablation.
\end{minipage}
\end{table}

\begin{figure}[htbp]
\centering
\includegraphics[width=\textwidth]{initialization_ablation.png}
\caption{Paired adversarial-training effects on exact-reserve mean squared
error. Negative values favor training. Bars are the mean paired squared-error
difference plus or minus 1.96 Monte Carlo standard errors across ten
replications, not confidence intervals for the reserve. Horizontal scales
differ across panels. These retrospective, approximate uncertainty summaries
are not a confirmatory dominance test.}
\label{fig:initialization}
\end{figure}
